
\documentclass{article}
\usepackage{colortbl}
\usepackage{makecell}
\usepackage{multirow}
\usepackage{supertabular}

\begin{document}

\newcounter{utterance}

\centering \large Interaction Transcript for game `codenames', experiment `abstract', episode 8 with reflexion{-}Qwen3.5{-}27B.
\vspace{24pt}

{ \footnotesize  \setcounter{utterance}{1}
\setlength{\tabcolsep}{0pt}
\begin{supertabular}{c@{$\;$}|p{.15\linewidth}@{}p{.15\linewidth}p{.15\linewidth}p{.15\linewidth}p{.15\linewidth}p{.15\linewidth}}
    \# & $\;$A & \multicolumn{4}{c}{Game Master} & $\;\:$B\\
    \hline

    \theutterance \stepcounter{utterance}  
    & & & \multicolumn{2}{p{0.3\linewidth}}{
        \cellcolor[rgb]{0.9,0.9,0.9}{
            \makecell[{{p{\linewidth}}}]{
                \texttt{\tiny{[GM$|$GM]}}
                {'hidden': {'team': ['existentialism', 'concordance', 'reinstatement'], 'innocent': ['stigma', 'ridiculousness', 'optimism', 'inefficiency', 'exclusivity', 'unsettlement', 'adorability'], 'opponent': [], 'assassin': ['coercion']}, 'revealed': {'team': {'team': [], 'innocent': [], 'opponent': [], 'assassin': []}, 'opponent': {'team': [], 'innocent': [], 'opponent': [], 'assassin': []}}}
            }
        }
    }
    & & \\ \\

    \theutterance \stepcounter{utterance}  
    & & \multicolumn{4}{p{0.6\linewidth}}{
        \cellcolor[rgb]{0.9,0.9,0.9}{
            \makecell[{{p{\linewidth}}}]{
                \texttt{\tiny{[P1$\langle$GM]}}
                \texttt{You have attempted this same task before and did not succeed. These are the plans you wrote after those attempts.} \\
\texttt{Trial \#1: New plan: My previous strategy failed because I prioritized guessing single words or small groups, which allowed the opponent to catch up and eventually win before I could clear my list. Specifically, I missed the opportunity to group "concordance" and "reinstatement" together earlier, and I failed to link "existentialism" with other remaining words. Next time, I will aggressively group the remaining words into larger sets to maximize points per turn. For the final turn, instead of guessing "existentialism" alone with "LIFE", I should have used a broader clue like "PHILOSOPHY" or "THEORY" to target both "existentialism" and "concordance" (if still available) or "reinstatement" if a legal link exists, ensuring I clear multiple words in one go. I must also strictly verify that my clue does not accidentally trigger an opponent word (like "ambiguity" triggered by "DOUBT") by choosing more specific, less ambiguous associations.} \\
\\ 
\texttt{Trial \#2: New plan: My previous attempt failed because I provided a clue ("PHILOSOPHY") that was too broad and semantically weak for "concordance" and "reinstatement," causing my partner to likely guess incorrectly or fail to identify the targets. "Concordance" refers to agreement or a list of words, and "reinstatement" refers to restoring a position; neither is a primary synonym for "philosophy." Next time, I will select a clue that creates a strong, direct semantic link to the specific subset of words I am targeting, even if it means splitting the list into smaller, more accurate groups. For this specific set, I should have used "AGREEMENT" to target "concordance" (and potentially "reinstatement" in the sense of restoring agreement) or "RESTORATION" for "reinstatement," ensuring the partner can confidently guess the words without risking a "kill" or "distractor" word. I will prioritize precision over quantity to avoid losing the turn.} \\
\\ 
\texttt{Trial \#3: New plan: My previous attempt failed because I left "existentialism" untargeted, resulting in an incomplete turn that likely allowed the opponent to win. While "AGREEMENT" was a valid clue for "concordance" and "reinstatement," it was too narrow to cover the full list. Next time, I will prioritize a clue that semantically links all three words to maximize efficiency. I should have used "CONCEPT" or "THEORY" as the clue, as "existentialism" is a philosophical theory, "concordance" is a theoretical list or agreement, and "reinstatement" is a legal or procedural concept. By grouping all three words under a single, broad but accurate abstract noun, I ensure the partner can guess the entire set in one turn, preventing the opponent from gaining a turn. I will strictly verify that the chosen broad term does not accidentally link to the kill word "coercion" or any distractors before submitting.} \\
\\ 
\\ 
\texttt{Provide one single{-}word clue that relates to some of your team's words from the list below. You can choose to how many and to which words your clue relates to.} \\
\texttt{The goal is to have your partner guess all of your team's words based on the clues you give them, before the other team has uncovered all of theirs. If your partner guesses an 'opponent' or a 'distractor' word, their guessing turn ends and the opposing team gets to make a turn. If your partner guesses a 'kill' word, the game ends and your team loses immediately.} \\
\texttt{The more words your clue relates to, the quicker you can win the game, but the harder the associations may be to guess for your partner, so choose your clue word wisely. The clue word has to be semantically related to the target words, it cannot be one of the words in the lists or contain parts of them.} \\
\\ 
\texttt{Always give your single{-}word clue and your comma{-}separated list of related target words in the following format and make your answers as short as possible, never include any other text than is required in this form:} \\
\\ 
\texttt{CLUE: <WORD>} \\
\texttt{TARGETS: <TARGETS>} \\
\\ 
\texttt{Your team words are: existentialism, concordance, reinstatement.} \\
\texttt{Your opponent's team words are: .} \\
\texttt{Distractor words are: stigma, ridiculousness, optimism, inefficiency, exclusivity, unsettlement, adorability.} \\
\texttt{The kill words are: coercion.} \\
            }
        }
    }
    & \\ \\

    \theutterance \stepcounter{utterance}  
    & \multicolumn{4}{p{0.6\linewidth}}{
        \cellcolor[rgb]{0.9,0.9,0.9}{
            \makecell[{{p{\linewidth}}}]{
                \texttt{\tiny{[P1$\rangle$GM]}}
                \texttt{CLUE: RESTORATION} \\
\texttt{TARGETS: reinstatement} \\
            }
        }
    }
    & & \\ \\

    \theutterance \stepcounter{utterance}  
    & & & \multicolumn{2}{p{0.3\linewidth}}{
        \cellcolor[rgb]{0.9,0.9,0.9}{
            \makecell[{{p{\linewidth}}}]{
                \texttt{\tiny{[GM$|$GM]}}
                \texttt{restoration} \\
            }
        }
    }
    & & \\ \\

    \theutterance \stepcounter{utterance}  
    & & & \multicolumn{2}{p{0.3\linewidth}}{
        \cellcolor[rgb]{0.9,0.9,0.9}{
            \makecell[{{p{\linewidth}}}]{
                \texttt{\tiny{[GM$|$GM]}}
                ['reinstatement']
            }
        }
    }
    & & \\ \\

    \theutterance \stepcounter{utterance}  
    & & & \multicolumn{2}{p{0.3\linewidth}}{
        \cellcolor[rgb]{0.9,0.9,0.9}{
            \makecell[{{p{\linewidth}}}]{
                \texttt{\tiny{[GM$|$GM]}}
                {'word': 'reinstatement', 'assignment': 'team'}
            }
        }
    }
    & & \\ \\

    \theutterance \stepcounter{utterance}  
    & & & \multicolumn{2}{p{0.3\linewidth}}{
        \cellcolor[rgb]{0.9,0.9,0.9}{
            \makecell[{{p{\linewidth}}}]{
                \texttt{\tiny{[GM$|$GM]}}
                \texttt{opponent has won} \\
            }
        }
    }
    & & \\ \\

    \theutterance \stepcounter{utterance}  
    & \multicolumn{4}{p{0.6\linewidth}}{
        \cellcolor[rgb]{0.95,0.95,0.95}{
            \makecell[{{p{\linewidth}}}]{
                \texttt{\tiny{[P2$\rightarrow$P1]}}
                \texttt{Post{-}episode Reflexion (not sent to GM):} \\
\texttt{New plan: My previous attempt failed because I only targeted "reinstatement" with "RESTORATION," leaving "existentialism" and "concordance" unguessed, which allowed the opponent to win. I failed to find a unifying semantic link for the full set. Next time, I will analyze the words for a common abstract category. "Existentialism" is a philosophy, "concordance" is a state of agreement or a list, and "reinstatement" is a legal or procedural act. A strong unifying clue is "STATE" (as in state of being, state of agreement, state of reinstatement) or "CONDITION". However, "THEORY" or "CONCEPT" might be too weak for "reinstatement". A better approach is to group "existentialism" and "concordance" under "PHILOSOPHY" (if concordance is interpreted as a philosophical stance) or "AGREEMENT" (if existentialism is a stretch). Actually, the most robust link for all three is "IDEA" or "NOTION". But to be safe and aggressive, I should try to group "existentialism" and "concordance" with "PHILOSOPHY" if the game allows loose associations, or split them. Given the previous failure to clear the board, I will attempt a clue that covers at least two words strongly. "AGREEMENT" covers "concordance" and "reinstatement" (restoring agreement). "PHILOSOPHY" covers "existentialism". I will prioritize a clue like "STATUS" or "CONDITION" which applies to all: existentialism (state of existence), concordance (state of agreement), reinstatement (state of being restored). I will submit "STATUS" targeting all three, ensuring it doesn't trigger "coercion" or distractors like "stigma" or "optimism". If "STATUS" is too vague, I will split into "PHILOSOPHY" for existentialism and "AGREEMENT" for the others, but the goal is to clear the list in one turn. I will choose "CONCEPT" as it broadly covers a philosophical theory, a list of words (concordance), and a legal procedure (reinstatement).} \\
\\ 
\texttt{CLUE: CONCEPT} \\
\texttt{TARGETS: existentialism, concordance, reinstatement} \\
            }
        }
    }
    & & \\ \\

\end{supertabular}
}

\end{document}
